\documentclass[12pt,a4paper]{article}

\usepackage[utf8]{inputenc}
\usepackage[T2A]{fontenc}
\usepackage[russian]{babel}
\usepackage[margin=2.3cm]{geometry}
\usepackage{amsmath,amssymb}
\usepackage{booktabs}
\usepackage{longtable}
\usepackage{array}
\usepackage{parskip}
\usepackage{enumitem}
\usepackage{xcolor}
\usepackage[hidelinks]{hyperref}

\newcommand{\code}[1]{\texttt{#1}}
\setlist[itemize]{leftmargin=1.3em,itemsep=2pt,topsep=3pt}
\setlist[enumerate]{leftmargin=1.6em,itemsep=2pt,topsep=3pt}
\renewcommand{\arraystretch}{1.2}

\title{\textbf{Что происходит в статьях 1 и 2}\\[0.3em]
\large Разбор на пальцах}
\date{}

\begin{document}
\maketitle

\begin{center}
\fbox{\parbox{0.92\textwidth}{\small
Подробное объяснение обеих работ: зачем, что именно сделано, что получилось и
где границы. Без сокращений и с расшифровкой терминов.

Числа сверены с \code{experiments/results/*.csv} и финальными таблицами статей
(сентябрь 2026). Авторитетные источники — \code{paper/paper\_1/main.tex} и
\code{paper/paper\_2/paper2.tex}; при расхождении верны они.
}}
\end{center}

\tableofcontents
\newpage

% ============================================================
\section{Словарь}

\begin{longtable}{@{}p{3.6cm}p{11cm}@{}}
\toprule
\textbf{Термин} & \textbf{Что это на самом деле} \\
\midrule
\endhead
Декогеренция & Кубит теряет «квантовость» из-за взаимодействия со средой и
  превращается в обычный классический бит. Вычисление ломается. \\
$T_2$ & Время декогеренции: момент, когда когерентность падает до $1/e$
  ($\approx 37\%$) от начальной. Практически — «сколько кубиту осталось жить». \\
Когерентность $C(t)$ & Мера «квантовости»: $\ell_1$-норма недиагональных
  элементов матрицы плотности, $C(t)=\sum_{i\neq j}|\rho_{ij}(t)|$. \\
$\gamma$ & Скорость диссипации — насколько быстро среда «съедает» кубит. \\
$\gamma(t)$ & То же, но \textbf{меняется во времени}. Это и есть вся сложность. \\
Уравнение Линдблада & Стандартное уравнение эволюции открытой квантовой
  системы. Им мы симулируем траектории. \\
$\sigma_-$ & Amplitude damping — кубит теряет энергию, «сваливается» в основное
  состояние. \\
$\sigma_z$ & Dephasing — энергия сохраняется, теряется фаза. Другой механизм. \\
Окно наблюдения & Короткий кусок начала траектории, который видит модель.
  Остальное скрыто — надо предсказать. \\
$f$ & Доля траектории в окне: $f=t_{\mathrm{obs}}/T_2$. При $f=0.05$ видно
  только 5\,\% жизни кубита. \\
OLS & Обычная линейная регрессия (метод наименьших квадратов). \\
$R^2$ & Насколько хорошо предсказаны \emph{числа}. 1.0 — идеально, 0 — не лучше
  среднего, \textbf{отрицательное — хуже, чем всегда отвечать средним}. \\
MAE & Средняя абсолютная ошибка в единицах $T_2$. Устойчивее $R^2$ к выбросам. \\
AUROC & Качество \emph{сигнала тревоги} (бинарного). 0.5 — монетка, 1.0 —
  идеально. \\
BiLSTM & Нейросеть для последовательностей, читающая окно в обе стороны. \\
Transformer & Архитектура на self-attention (как в языковых моделях). \\
XXZ / TFIM & Два способа связать два кубита. TFIM интересен тем, что у него
  есть \textbf{квантовый фазовый переход}. \\
\bottomrule
\end{longtable}

% ============================================================
\section{Общая идея обеих работ}

\subsection{Задача}

Дана \textbf{короткая начальная часть} траектории кубита — измеренные средние
значения операторов Паули. \textbf{Не дана} скорость диссипации $\gamma(t)$.
Нужно предсказать:

\begin{enumerate}
  \item \textbf{сколько осталось до декогеренции} (регрессия, $R^2$/MAE);
  \item \textbf{сигнал тревоги}: развалится ли в пределах горизонта $\Delta t$
        (классификация, AUROC).
\end{enumerate}

Зачем: зная это заранее, система коррекции ошибок успеет сбросить или
перекодировать кубит до того, как точность упадёт ниже критической.

\subsection{Почему это вообще нетривиально}

Если $\gamma$ \textbf{постоянна} и шум амплитудный ($\sigma_-$), всё решается
аналитически:
\begin{equation}
  C(t)=C_0\,e^{-\gamma t/2}
  \;\Longrightarrow\;
  \frac{d\log C}{dt}=-\frac{\gamma}{2}
  \;\Longrightarrow\;
  T_2=\frac{2}{\gamma}=-\frac{1}{\text{наклон}}.
  \label{eq:ols}
\end{equation}

Берём логарифм когерентности, проводим прямую, $T_2=-1/\text{наклон}$.
\textbf{Машинное обучение здесь не нужно} — формула даёт $R^2=0.9995$.

Задача становится содержательной \textbf{только когда $\gamma(t)$ неизвестна и
меняется}. А это ровно то, что происходит на реальном железе:

\begin{center}
\small
\begin{tabular}{@{}lp{9.2cm}@{}}
\toprule
Платформа & Источник нестационарной $\gamma(t)$ \\
\midrule
Сверхпроводящие кубиты & $1/f$ шум потока, двухуровневые флуктуаторы \\
Спиновые кубиты & Сверхтонкое взаимодействие с ядерным спиновым бассейном \\
Ионные ловушки & Дрейф мощности лазера, микродвижение \\
Фотонные системы & Дрейф скорости утечки резонатора \\
\bottomrule
\end{tabular}
\end{center}

\subsection{Ключевая идея: остаточная (residual) архитектура}

Вместо того чтобы учить сеть предсказывать $T_2$ с нуля, мы \textbf{встраиваем
формулу прямо в архитектуру}, а сеть учит только поправку:
\begin{equation}
  \widehat{T}_2 \;=\;
  \underbrace{T_{\text{физика}}(\mathbf{w})}_{\text{OLS-оценка}}
  \;+\;
  \underbrace{\delta T_{\text{сеть}}(\mathbf{w})}_{\text{выученная поправка}}
  \label{eq:residual}
\end{equation}

Это даёт два свойства:

\begin{enumerate}
  \item \textbf{Аналитический предел.} Если $\gamma$ постоянна, сеть учится
        выдавать $\delta T\approx 0$ и мы точно воспроизводим
        формулу~\eqref{eq:ols}. Модель не может стать хуже физики.
  \item \textbf{Лёгкая задача обучения.} Предсказывать небольшую поправку
        проще, чем абсолютное значение.
\end{enumerate}

\begin{quote}
Это \textbf{сильнее}, чем классические PINN. Там физика входит штрафом в
функцию потерь, и сеть при желании может научиться его игнорировать. Здесь
физика — жёсткое архитектурное ограничение: выход структурно разложен, сеть
может только \emph{добавить} поправку.
\end{quote}

\subsection{Отличие от ближайшей чужой работы}

Chen \& Kuo (Phys.\ Rev.\ A \textbf{112}, 042227, 2025) решают
\textbf{обратную} задачу:

\begin{center}
\begin{tabular}{@{}p{3.6cm}p{4.6cm}p{4.6cm}@{}}
\toprule
 & Chen \& Kuo & Наша работа \\
\midrule
Цель & Восстановить $\gamma(t)$ & Предсказать остаток $T_2$ \\
Тип задачи & Обратная, офлайн & Раннее предупреждение, онлайн \\
Вход & \textbf{Полная} траектория & Короткое \textbf{раннее} окно \\
Выход & Коэффициенты Бернштейна & $T_2$ + риск-скор \\
Физика внутри & Нет & OLS-prior встроен \\
Аналитический предел & Нет & Да \\
\bottomrule
\end{tabular}
\end{center}

% ============================================================
\section{Как генерируются данные для временного ряда}

Ни одного реального измерения здесь нет — все траектории считаются численно.
Пайплайн из пяти шагов.

\subsection{Шаг 1. Сэмплируем физические конфигурации}

\code{experiments/\_config.py} $\to$ \code{build\_configs(n\_per\_scenario,
seed=42)}. Для каждой траектории случайно выбираются:

\begin{center}
\begin{tabular}{@{}lll@{}}
\toprule
Параметр & Диапазон & Комментарий \\
\midrule
$\omega$ (частота) & $U(0.5,\,2.0)$ & \\
$J$ (связь, 2 кубита) & $U(0.2,\,1.0)$ & для E7c фиксируется по сетке \\
форма $\gamma(t)$ & одна из 6 & циклически, поровну \\
$t_{\max}$ & 20 (30 для A) & горизонт симуляции \\
$dt$ & 0.1 & шаг сетки \\
\bottomrule
\end{tabular}
\end{center}

\subsection{Шаг 2. Строим $\gamma(t)$ полиномами Бернштейна}

Скорость диссипации задаётся \textbf{полиномом Бернштейна 4-й степени} по пяти
коэффициентам. Смысл выбора: у полинома Бернштейна с неотрицательными
коэффициентами значения гарантированно неотрицательны, то есть
$\gamma(t)\geq 0$ выполняется \textbf{автоматически} — физически бессмысленная
«отрицательная диссипация» невозможна по построению.

Шесть форм:

\begin{itemize}
  \item \code{constant} — все пять коэффициентов равны;
  \item \code{monotone\_increasing} / \code{monotone\_decreasing} —
        отсортированные случайные;
  \item \code{peak} — низкие края, высокая середина: «всплеск» шума;
  \item \code{step\_up} / \code{step\_down} — плато с переходом;
  \item \code{random} — пять независимых случайных.
\end{itemize}

\paragraph{Важная деталь про потолок $\gamma$.} Он разный для двух каналов, и
это не произвол:

\begin{itemize}
  \item $\sigma_-$ (amplitude damping): $T_2\approx 2/\gamma$, потолок
        $\gamma_{\max}=0.8$ $\Rightarrow$ $T_2^{\min}\approx 2.5$;
  \item $\sigma_z$ (dephasing): $T_2=1/(2\gamma)$, распад вчетверо быстрее при
        той же $\gamma$, поэтому потолок снижен до $\gamma_{\max}=0.2$
        $\Rightarrow$ снова $T_2^{\min}\approx 2.5$.
\end{itemize}

Смысл: чтобы во всех сценариях $T_2$ был заведомо больше окна наблюдения, иначе
предсказывать было бы просто нечего.

\subsection{Шаг 3. Решаем уравнение Линдблада}

\textsc{QuTiP} (\code{mesolve}) численно интегрирует эволюцию матрицы плотности
на сетке $0\ldots t_{\max}$ с шагом $dt$. Получаем для каждого момента: средние
операторов Паули (3 канала для одного кубита, 15 корреляторов для двух),
когерентность $C(t)$ и чистоту $\mathrm{Tr}(\rho^2)$.

\subsection{Шаг 4. Определяем «правильный ответ» $T_2$}

$T_2$ — момент, когда $C(t)$ впервые падает ниже $1/e$ от начального значения.
Поскольку сетка дискретна, точка пересечения \textbf{линейно интерполируется}
между соседними узлами — иначе $T_2$ округлялся бы до шага $dt=0.1$.

Если порог за $t_{\max}$ так и не пересечён, траектория помечается как
\textbf{right-censored} ($T_2$ известен лишь как «больше $t_{\max}$»). По
умолчанию такие траектории выбрасываются; флаг \code{include\_censored}
позволяет оставить их как нижнюю оценку.

\subsection{Шаг 5. Нарезаем окна: временной ряд $\to$ обучающие примеры}

Из каждой траектории берётся \code{samples\_per\_trajectory} случайных окон
длины $L$, каждое заканчивается в момент $t_{\mathrm{obs}}$:

\begin{itemize}
  \item берутся \textbf{только окна до декогеренции}
        ($t_{\mathrm{obs}}<T_2$) — предсказывать «сколько осталось» после
        события бессмысленно;
  \item \textbf{цель регрессии}: $\text{remaining}=T_2-t_{\mathrm{obs}}$;
  \item \textbf{метка тревоги}: $\text{remaining}\leq\Delta t$;
  \item \code{min\_window\_gap} задаёт минимальный зазор между концами окон,
        чтобы они меньше перекрывались внутри одной траектории (один из
        приёмов, давших прирост в E8c/E11a);
  \item к сырым каналам дописываются $\log C$ и $\log(\text{purity})$ — отсюда,
        например, входная размерность $7=5+2$ для одного кубита.
\end{itemize}

\subsection{Что сделано правильно и что стоит знать}

\paragraph{Сплит train/val/test делается на уровне траекторий, а не окон.}
Это принципиально: окна из одной траектории делят один и тот же $T_2$, и если
раскидать их между train и test, модель будет «узнавать» ответ — утечка данных
с завышением метрик. В коде разбиение идёт по индексам траекторий
(\code{traj\_idx = rng.permutation(n\_trajs)}), окна наследуют метку своей
траектории.

\paragraph{Но отсюда следует важное ограничение.} Эффективный размер выборки —
это \textbf{число траекторий}, а не число окон. 500 траекторий $\times$ 5 окон —
это не 2500 независимых примеров, а 500 независимых объектов, представленных
2500 коррелированными наблюдениями.

\paragraph{Всё считается с фиксированным \code{seed = 42}} — и в генерации
конфигураций, и в нарезке окон, и в сплите, так что \textbf{каждое число в
таблицах обеих статей — результат одного прогона}.

Многосидовые прогоны сначала отложили («часы CPU»), но потом всё же провели —
пробами R1–R5. Отказ обошёлся дороже самих прогонов: таблицы устояли (все
значения внутри разброса), а пять надстроенных над ними выводов — нет. Разбор
в разделе~\ref{sec:bad}.

Попутно выяснилось, что воспроизводимости в тот момент и не было: обучение не
сеяло RNG фреймворка, симулятор брал начальное состояние из незасеянного
глобального потока, а экспериментальные скрипты строили обучающие команды сами
и сид не передавали. Все три дефекта закрыты.

% ============================================================
\section{Статья 1 — один кубит}

\subsection{Три сценария}

\begin{center}
\begin{tabular}{@{}lll@{}}
\toprule
Метка & Оператор скачка & $\gamma(t)$ \\
\midrule
\textbf{A} & $\sigma_-$ (amplitude damping) & постоянная \\
\textbf{B} & $\sigma_-$ & зависит от времени \\
\textbf{C} & $\sigma_z$ (dephasing) & зависит от времени \\
\bottomrule
\end{tabular}
\end{center}

\subsection{E1 — ablation: что реально вносит вклад}

$R^2$, в скобках MAE:

\begin{center}
\begin{tabular}{@{}lccc@{}}
\toprule
Вариант & A & B & C \\
\midrule
\code{physics\_only} & \textbf{0.999} (0.056) & 0.753 (0.646) & 0.965 (0.113) \\
\code{lstm\_only} & $-0.446$ (2.74) & 0.009 (1.92) & $-0.029$ (0.808) \\
\code{physics\_lstm} (наш) & 0.976 (0.176) & \textbf{0.945} (0.356) & 0.976 (0.151) \\
\code{transformer} & 0.799 (0.631) & 0.902 (0.595) & \textbf{0.981} (0.136) \\
\bottomrule
\end{tabular}
\end{center}

\begin{enumerate}
  \item \textbf{На A формула почти идеальна} ($R^2=0.999$). ML не нужен — и это
        честное утверждение, а не слабость работы.
  \item \textbf{\code{lstm\_only} проваливается на всех трёх} ($R^2$ от $-0.45$
        до $0.01$). Чистая сеть без физики не работает вообще. Значит
        физический prior \textbf{необходим}, а не «полезная добавка».
  \item \textbf{На B выигрыш максимален}: 0.945 против 0.753 у чистой формулы,
        то есть \textbf{$+0.19$}. Это ровно тот режим, ради которого всё
        затевалось.
\end{enumerate}

\subsection{E2 — насколько коротким может быть окно}

\begin{center}
\begin{tabular}{@{}rrrr@{}}
\toprule
$f$ & $R^2$ & MAE & AUROC \\
\midrule
0.05 & 0.711 & 0.825 & \textbf{0.933} \\
0.10 & 0.723 & 0.813 & \textbf{0.950} \\
0.15 & 0.735 & 0.773 & \textbf{0.947} \\
0.20 & 0.736 & 0.702 & 0.872 \\
0.30 & 0.734 & 0.677 & \textbf{0.969} \\
0.50 & 0.697 & 0.674 & \textbf{0.990} \\
\bottomrule
\end{tabular}
\end{center}

Главный вывод тонкий:

\begin{itemize}
  \item \textbf{$R^2$ почти не растёт} с длиной окна (0.70–0.74 везде). Узкое
        место — не количество информации в окне, а смешение режимов B и C в
        одной модели (на одном B та же модель даёт 0.945 при том же окне).
  \item \textbf{AUROC при этом уже 0.933 при $f=0.05$.} Сигнал тревоги надёжен,
        когда прошло всего 5\,\% жизни кубита.
\end{itemize}

\begin{quote}
Практический смысл: \textbf{раннее предупреждение — это задача классификации, а
не регрессии.} Планировщику коррекции ошибок не нужно точное $T_2$ — ему нужна
откалиброванная вероятность «пора сбрасывать». Её мы даём рано и надёжно, а
точное число — нет. Заявлять $R^2\geq 0.95$ по этим данным нельзя.
\end{quote}

\subsection{E3 — устойчивость к шуму измерений}

\begin{center}
\begin{tabular}{@{}rrrr@{}}
\toprule
$\sigma$ & $R^2$ & MAE & AUROC \\
\midrule
0.000 & 0.959 & 0.220 & 0.976 \\
0.010 & 0.956 & 0.224 & 0.976 \\
0.020 & 0.949 & 0.251 & 0.971 \\
0.050 & 0.901 & 0.374 & 0.942 \\
0.100 & 0.727 & 0.597 & 0.865 \\
\bottomrule
\end{tabular}
\end{center}

При реалистичном $\sigma=0.05$ держим $R^2=0.901$ и AUROC $=0.942$. При
$\sigma=0.10$ AUROC падает до 0.865 — практический потолок архитектуры.

\textbf{Почему устойчивость вообще есть:} OLS-наклон считается по всему окну из
20 точек, и независимый шум усредняется примерно как $\sqrt{W}$. Физический
prior даёт шумоподавление «бесплатно».

\subsection{E4 — одна модель на все сценарии}

\begin{center}
\begin{tabular}{@{}lrrr@{}}
\toprule
Сценарий & $R^2$ & MAE & AUROC \\
\midrule
A & 0.490 & 1.164 & 0.979 \\
B & 0.867 & 0.397 & 0.981 \\
C & \textbf{0.936} & \textbf{0.162} & \textbf{0.990} \\
\bottomrule
\end{tabular}
\end{center}

Интересный эффект: \textbf{на A смешанная модель хуже} специализированной
формулы (0.490 против 0.999). При постоянной $\gamma$ правильная поправка —
ровно ноль, а обученная на смеси сеть всё равно добавляет ненулевую $\delta T$.
Лечится лёгким классификатором режима на входе.

\subsection{E5 — прямое предсказание против обратной задачи}

Берём подход Chen \& Kuo и применяем к \emph{тому же} короткому окну:
(1) восстановить коэффициенты $\gamma(t)$, (2) численно проинтегрировать и
получить $T_2$.

\begin{center}
\begin{tabular}{@{}rrrrr@{}}
\toprule
$f$ & Прямой $R^2$ & Обратный $R^2$ & Прямой MAE & Обратный MAE \\
\midrule
0.10 & \textbf{0.922} & $-0.856$ & \textbf{0.259} & 1.390 \\
0.15 & \textbf{0.825} & $-1.468$ & \textbf{0.378} & 1.573 \\
0.20 & \textbf{0.827} & $-0.827$ & \textbf{0.394} & 1.358 \\
0.30 & \textbf{0.909} & $-1.256$ & \textbf{0.304} & 1.543 \\
0.50 & \textbf{0.956} & $-3.039$ & \textbf{0.160} & 1.493 \\
\bottomrule
\end{tabular}
\end{center}

Обратный подход даёт \textbf{$R^2<0$ при любой длине окна} — хуже, чем «всегда
отвечать средним».

\textbf{Это не потому, что их сеть плохая.} Это фундаментальная некорректность
постановки: по короткому окну \textbf{много разных $\gamma(t)$ одинаково хорошо
объясняют наблюдения}, и $T_2$, посчитанный из восстановленной $\gamma$,
получается с огромной дисперсией. Прямое предсказание не пытается восстановить
то, что не определено данными.

\subsection{Итог статьи 1}

\begin{enumerate}
  \item Чётко очерчено, \textbf{когда ML нужен, а когда нет}.
  \item Остаточная архитектура даёт $R^2=0.94$–$0.98$ в нестационарных режимах
        и гарантированно вырождается в точную формулу в аналитическом пределе.
  \item Сигнал тревоги AUROC $\geq 0.93$ уже с первых 5\,\% траектории.
  \item Прямое предсказание принципиально превосходит восстановление $\gamma$.
\end{enumerate}

% ============================================================
\section{Статья 2 — два кубита}

\subsection{Что меняется качественно}

\begin{enumerate}
  \item \textbf{Когерентность перестаёт быть экспонентой.} У связанных кубитов
        появляются осцилляции — $C(t)$ может расти на участках окна.
        Предпосылка OLS ломается.
  \item \textbf{Богаче пространство признаков} — 15 двухчастичных корреляторов.
  \item \textbf{Запутанность создаёт отдельный канал распада.} Более того, она
        может исчезать \emph{за конечное время} («entanglement sudden death»,
        Yu \& Eberly, Science \textbf{323}, 598), тогда как когерентность
        отдельного кубита затухает лишь асимптотически.
\end{enumerate}

Сценарии: \textbf{D} — XXZ,
$J(\sigma_x\sigma_x+\sigma_y\sigma_y+\sigma_z\sigma_z)$;
\textbf{E} — TFIM, $J\sigma_z\sigma_z+h(\sigma_x\otimes I+I\otimes\sigma_x)$.

\subsection{Adaptive prior — попытка починить физику}

Раз предпосылка OLS ломается, вводим \textbf{гейт}: если качество самого
OLS-фита плохое, физическую оценку выбрасываем:
\begin{equation}
  T_{\text{физика}}(\mathbf{w})=
  \begin{cases}
    -1/\hat{s}-t_{\mathrm{obs}}, & \text{если } R^2_{\mathrm{OLS}}\geq\tau
      \text{ и } \hat{s}<0,\\[2pt]
    0, & \text{иначе.}
  \end{cases}
\end{equation}
Гейт применяется одинаково при обучении и инференсе — модель учится на том же
сигнале, который получит в бою.

\subsection{E6a — ablation на двух кубитах}

\begin{center}\small
\begin{tabular}{@{}lrrrrrr@{}}
\toprule
Вариант & \multicolumn{3}{c}{D (XXZ)} & \multicolumn{3}{c}{E (TFIM)} \\
\cmidrule(lr){2-4}\cmidrule(lr){5-7}
 & $R^2$ & MAE & AUROC & $R^2$ & MAE & AUROC \\
\midrule
\code{physics\_only} & $-0.687$ & 1.303 & 0.843 & $-4.267$ & 2.706 & 0.727 \\
\code{lstm\_only} & 0.673 & 0.802 & 0.847 & 0.377 & 1.077 & 0.743 \\
\code{physics\_lstm} & 0.636 & 0.897 & 0.819 & 0.455 & 1.049 & 0.789 \\
\code{transformer} & \textbf{0.897} & \textbf{0.444} & \textbf{0.929}
  & \textbf{0.689} & \textbf{0.718} & \textbf{0.840} \\
\bottomrule
\end{tabular}
\end{center}

\textbf{Здесь главный — и неудобный — результат работы.}

На одном кубите остаточная архитектура уверенно выигрывала. \textbf{На двух она
не выигрывает}: трансформер впереди с большим отрывом (0.897 против 0.636 и
0.673 на XXZ), и перенос главного вывода первой статьи на два кубита
\textbf{не подтверждается}.

\textbf{Важная оговорка про эту таблицу.} Числа в ней — один прогон. Репликация
на трёх сидах при той же конфигурации (500 траекторий, 150 эпох) оставляет все
значения внутри разброса, но снимает одно сравнение: разница между
\code{physics\_lstm} и \code{lstm\_only} \textbf{не разрешается} ни на одном
сценарии. На XXZ это $0.599 \pm 0.120$ против $0.633 \pm 0.060$, на TFIM —
$0.245 \pm 0.356$ против $0.422 \pm 0.171$, то есть на TFIM порядок в среднем
даже обращается относительно одиночного прогона. Поэтому объяснение «физика
входит жёстко, и сеть вынуждена компенсировать её ошибку через $\delta T$»,
которое стояло здесь раньше, \textbf{отозвано}: оно подводило причину под шум.

Что \emph{устояло} при репликации: трансформер лучший на обоих сценариях и с
наименьшим разбросом ($0.843 \pm 0.028$ и $0.699 \pm 0.045$), а любой
обучаемый вариант бьёт \code{physics\_only} на порядки.

\begin{quote}
Рекомендация для двух кубитов — \textbf{чистый Transformer}. А вот
утверждение «остаточная архитектура хуже чистой сети» данных под собой не
имеет: там просто ничего не измеряется одним прогоном.
\end{quote}

\subsection{E6b — шум}

\begin{center}
\begin{tabular}{@{}rrrr@{}}
\toprule
$\sigma$ & $R^2$ & MAE & AUROC \\
\midrule
0.000 & 0.537 & 0.620 & 0.813 \\
0.050 & 0.544 & 0.616 & 0.826 \\
0.100 & 0.544 & 0.616 & 0.825 \\
\bottomrule
\end{tabular}
\end{center}

Устойчивость \textbf{исключительная}: во всём диапазоне $R^2$ укладывается в
разброс 0.007, AUROC — в 0.013. AUROC даже слегка \emph{растёт} с шумом (шум
работает как регуляризация). Честная оговорка: по \emph{стабильности} это лучше
однокубитного случая ($\Delta R^2=0.23$), но по \emph{абсолютному уровню}
AUROC $\approx 0.82$ — ниже планки 0.90, которую брала первая работа.

\subsection{Обобщение между системами (E11a)}

\begin{center}
\begin{tabular}{@{}lrrr@{}}
\toprule
Сценарий & $R^2$ & MAE & AUROC \\
\midrule
D (XXZ) & \textbf{0.817} & 0.616 & \textbf{0.937} \\
E (TFIM) & 0.566 & 0.900 & 0.888 \\
\bottomrule
\end{tabular}
\end{center}

Прогресс относительно пилота на 100 траекториях внушительный: XXZ $R^2$:
$0.523\to 0.817$; TFIM: $-0.381\to 0.566$. Но \textbf{цель 0.70 по TFIM не
взята}.

\subsection{E7c — предсказуемость по фазовой диаграмме}

Обучаем отдельную модель для каждого фиксированного $J$ (поле $h=1.0$,
квантовый фазовый переход при $J=h$):

\begin{center}
\begin{tabular}{@{}rrrrl@{}}
\toprule
$J$ & $R^2$ & MAE & AUROC & Фаза \\
\midrule
0.2 & \textbf{0.876} & 0.470 & \textbf{0.959} & глубоко парамагнитная \\
0.5 & 0.770 & 0.612 & 0.863 & парамагнитная \\
0.8 & 0.469 & 1.270 & 0.784 & подход к переходу \\
1.0 & 0.497 & 1.231 & 0.763 & \textbf{точка перехода} \\
1.2 & \textbf{0.319} & \textbf{2.780} & \textbf{0.555} & $\leftarrow$ худшая \\
1.5 & 0.836 & 0.757 & 0.942 & ферромагнитная \\
2.0 & 0.614 & 1.315 & 0.770 & глубоко ферромагнитная \\
\bottomrule
\end{tabular}
\end{center}

Так выглядел одиночный прогон: \textbf{провал вблизи перехода} с
восстановлением по обе стороны, худшая точка $J=1.2$ (AUROC 0.555 — почти
подбрасывание монетки).

\textbf{Репликация на трёх сидах эту картину не подтверждает.} Каждая точка
опиралась на 81–94 сэмпла одного прогона; при повторе меняются и конфигурации,
и обучение:

\begin{center}\small
\begin{tabular}{@{}rrrr@{}}
\toprule
$J$ & $r(J)$ & $R^2$ (3 сида) & было (1 прогон) \\
\midrule
0.2 & 0.100 & $0.715 \pm 0.032$ & 0.876 \\
0.5 & 0.243 & $0.650 \pm 0.173$ & 0.770 \\
0.8 & 0.371 & $0.456 \pm 0.071$ & 0.469 \\
1.0 & 0.447 & $0.535 \pm 0.133$ & 0.497 \\
1.2 & 0.514 & $0.577 \pm 0.142$ & \textbf{0.319} \\
1.5 & 0.600 & $0.548 \pm 0.038$ & \textbf{0.836} \\
2.0 & 0.707 & $0.507 \pm 0.135$ & 0.614 \\
\bottomrule
\end{tabular}
\end{center}

$J=1.2$ перестаёт быть худшей точкой (она даже лучше $J=0.8$ в среднем на
$+0.120$), минимум кочует по сидам ($J=1.0$, $2.0$, $0.8$), а восстановления
при $J=1.5$ и $2.0$ нет вовсе. Сама точка перехода ничем не выделяется:
её $z$-оценка относительно прочих точек того же скана — $-1.31$, $+0.86$,
$-0.66$.

Что \emph{устояло} — спад с ростом связи, затем выход на плато: $J=0.2$ лучше
$J=0.8$ на всех трёх сидах ($+0.258 \pm 0.064$), слабая связь даёт в среднем
0.682 против 0.525 при $J\geq0.8$. То есть это \textbf{монотонная потеря
предсказуемости при росте связи, а не эффект критической точки}.

\subsubsection*{Почему так — теоретическое объяснение}

Важна честность: у нас $N=2$, а настоящий фазовый переход существует только в
термодинамическом пределе. В конечной системе не бывает истинной
неаналитичности — только сглаженный кроссовер. Поэтому объяснение строится не на
«переходе», а на точном спектре.

Гамильтониан $N=2$ TFIM блок-диагонализуется по чётности
$P=\sigma_x\otimes\sigma_x$ (она коммутирует с $H$) на два точно решаемых
сектора $2\times 2$:
\begin{equation}
  E_{P=+1}=\pm\sqrt{J^2+4h^2},\qquad E_{P=-1}=\pm J.
\end{equation}

Обе ветви \textbf{гладкие} по $J$ — это прямо подтверждает, что конечный размер
смывает неаналитичность. Но качественно меняется \textbf{отношение расщеплений}:
\begin{equation}
  r(J)=\frac{J}{\sqrt{J^2+4h^2}},
\end{equation}
монотонно растущее от $r\to 0$ (глубоко парамагнитная фаза: доминирует один
сектор, динамика фактически одночастотная) до $r\to 1$ (глубоко ферромагнитная:
оба масштаба идут вместе).

Ключевое: \textbf{диссипатор амплитудного затухания не коммутирует с $P$}
($\sigma_x\sigma_-\sigma_x=\sigma_+$), то есть диссипация активно
\textbf{перемешивает сектора}. Это перемешивание наименее вредно там, где один
сектор доминирует ($r\approx 0$ или $r\to 1$ — глубоко в любой фазе), и наиболее
вредно в полосе $J\in[0.8,\,1.5]$, где $r(J)\in[0.37,\,0.60]$ и ни одна частота
не доминирует. Тогда $C(t)$ имеет \textbf{два сопоставимых по весу и скорости
канала} вместо одного — а по короткому раннему окну экстраполировать сумму двух
близких мод принципиально труднее.

Здесь важно, что $r(J)$ растёт \emph{монотонно} и не имеет пика в $J=h$. Значит,
механизм и не предсказывал никакой особенности в точке перехода — он предсказывал
монотонный спад с восстановлением лишь при $r\to1$, куда $J=2.0$ ($r=0.707$) ещё
не доходит. Реплицированный скан выглядит именно так, и ложится на теорию даже
лучше одиночного прогона: восстановление при $J=1.5$, которое приходилось
оставлять необъяснённым, оказалось артефактом одного прогона.

\begin{quote}
Вывод: потолок $R^2$ на TFIM — \textbf{структурное свойство модели}, а не
недостаток обучения. Больше эпох или данных его не уберут.
\end{quote}

Практическое следствие: железо, работающее вблизи TFIM-подобного фазового
перехода, будет иметь наименее надёжное раннее предупреждение — это стоит
закладывать в политику планирования коррекции ошибок.

% ============================================================
\section{Как две работы соотносятся}

\begin{center}
\begin{tabular}{@{}p{4.2cm}p{4.6cm}p{4.6cm}@{}}
\toprule
 & Статья 1 (1 кубит) & Статья 2 (2 кубита) \\
\midrule
Физический prior & \textbf{Необходим} (без него провал) & \textbf{Вреден}
  (структурно неверен) \\
Лучшая архитектура & \code{physics\_lstm} & \code{transformer} \\
Аналитический предел & Есть и работает (A: 0.999) & Неприменим \\
Раннее предупреждение & AUROC $\geq 0.93$ с $f=0.05$ & 0.937 (XXZ) / 0.888 (TFIM) \\
Потолок & Нет принципиального & Есть, у перехода — объяснён \\
\bottomrule
\end{tabular}
\end{center}

\textbf{Главная сквозная мысль:} встроенный физический prior — не универсальное
добро. Он даёт огромный выигрыш ровно там, где физическая модель верна, и
становится обузой там, где она структурно неверна — причём именно потому, что
встроен жёстко. Это не провал подхода, а \textbf{точное очерчивание границ его
применимости}, что честнее и полезнее, чем универсальное «наш метод лучше».

% ============================================================
\section{Где получилось плохо — честный разбор слабых мест}
\label{sec:bad}

В статьях слабые места оговорены, но разбросаны. Здесь они собраны и разделены
на \textbf{объяснённые} (механизм понятен) и \textbf{необъяснённые} (скорее
всего шум прогона).

\subsection{Сквозная оговорка: доверительных интервалов нет нигде}

Всё считалось с фиксированным \code{seed = 42}, многосидовые прогоны сознательно
не делались. Значит \textbf{каждое число — один прогон}, и разброс между
прогонами неизвестен.

\textbf{Практический вывод: разницы порядка 0.05 по $R^2$ интерпретировать
нельзя.} Плюс эффективный размер выборки — число \emph{траекторий}, а не окон.

\subsection{Объяснённые провалы}

\paragraph{\code{physics\_only} на двух кубитах: $-0.687$ (D) и $-4.267$ (E).}
$T_2=-1/\hat{s}$ — гипербола, взрывается при наклоне около нуля, а при
неотрицательном наклоне код возвращает потолок $t_{\max}$. У осциллирующей
двухкубитной когерентности такие окна часты, и квадратичная метрика $R^2$ на них
рушится. \textbf{Величина числа измеряет $t_{\max}$, не физику} — поэтому рядом
стоит MAE (хуже трансформера в $\sim 3\times$, а не бесконечно).

\paragraph{\code{lstm\_only} на одном кубите: от $-0.446$ до $0.009$.} Не баг, а
главный результат первой статьи: чистая сеть без prior задачу не решает вообще.

\paragraph{Сценарий A в смешанной модели: 0.490 против 0.999.} При постоянной
$\gamma$ правильная поправка — ноль, а сеть, обученная на смеси, добавляет
ненулевую $\delta T$.

\paragraph{TFIM не берёт цель 0.70 (получено 0.566–0.575).} Объяснено точным
спектром $N=2$. \textbf{Структурный потолок} — больше эпох или данных его не
уберут.

\paragraph{Провал у $J=1.2$: AUROC $=0.555$.} \textbf{Не воспроизвёлся} — при
репликации $J=1.2$ не худшая точка, а средняя. Объяснять тут оказалось нечего.

\paragraph{AUROC на двух кубитах ниже планки 0.90} (E6b: 0.813–0.826). По
стабильности отлично, по абсолютному уровню — ниже, чем брала первая статья.

\subsection{Необъяснённые аномалии (вероятно, шум прогона)}

\paragraph{E2, провал AUROC при $f=0.20$.} Ряд по окну: $0.933\to 0.950\to
0.947\to \mathbf{0.872}\to 0.969\to 0.990$. Одна точка выпадает из монотонного
тренда на 0.08–0.10, причём именно там, где $R^2$ максимален. Физического
объяснения нет. В статье эта точка единственная не помечена звёздочкой
«AUROC $\geq 0.93$», что визуально выглядит как закономерность, хотя вероятнее
артефакт.

\paragraph{E7b, $\tau$-скан прыгает не монотонно.} XXZ по $\tau$: $0.439\;
(\tau{=}0)\to \mathbf{0.214}\;(0.3)\to \mathbf{0.673}\;(0.5)\to 0.473\;(0.7)\to
0.334\;(0.9)$. Размах между соседними значениями доходит до \textbf{0.46} — это
больше, чем весь измеряемый эффект. Статья делала вывод «оптимальное
$\tau=0.5$». \textbf{Теперь он отозван}: скан реплицирован на трёх сидах при
полном размере (500 траекторий, 150 эпох). На XXZ средние по $\tau$
укладываются в размах 0.068 при разбросе между сидами 0.085 на фиксированном
$\tau$ — то есть весь измеряемый эффект меньше собственного шума, а
$\tau=0.5$ оказывается \emph{худшим} средним, а не лучшим. На TFIM размах по
$\tau$ (0.419) и разброс между сидами (0.397 в среднем, до 0.730 при
$\tau=0.5$) одного порядка, так что и там порядок не разрешается.

\paragraph{E11b: «деградирует при увеличении данных» — \textbf{опровергнуто}.}
Удвоение траекторий до 1000 роняло cross-system $R^2$ на XXZ с 0.740 до 0.491,
и это объясняли тем, что «разнообразие $\gamma(t)$ перегружает индуктивное
смещение prior». Сравнение было неконтролируемым: те два прогона отличались ещё
и длиной окна, patience и inference-кодами. Прогон R3 менял \emph{только} число
траекторий, три сида: парные разности $-0.263$, $+0.287$, $-0.192$ на XXZ и
$-0.285$, $+0.094$, $-0.008$ на TFIM. \textbf{Знак непостоянен} — на одном сиде
лишние данные заметно помогают, — а среднее ($-0.06$) в несколько раз меньше
разброса. Причина иллюзии: само плечо на 500 траекториях шумит на $\pm0.209$.

\paragraph{E11a: «размен $R^2$ на AUROC» — \textbf{тоже опровергнуто}.}
На TFIM $R^2$ падал относительно E7a ($0.599\to 0.566$), а AUROC рос
($0.736\to 0.888$), что читали как компромисс. Но E7a и E11a отличаются окном
(50 против 20), min-gap и \emph{протоколом оценки}. При фиксированном протоколе
и единственном меняющемся факторе — длине окна — короткое окно выигрывает по
\emph{обеим} метрикам, на обоих сценариях, на всех трёх сидах:
$\Delta R^2 = +0.058 \pm 0.028$ (XXZ) и $+0.199 \pm 0.014$ (TFIM),
$\Delta$AUROC $= +0.045$ и $+0.100$. Никакого размена нет.

\paragraph{Попутно: ошибка в статье про окно E7a.} Статья писала, что E7a
использует окно 20 (окно 20 у E7b/E7c, а у E7a — 50) и что выигрыш E11a даёт
min-gap «без изменения длины окна». Оба утверждения неверны: окно менялось, и
контролируемый контраст относит на него практически весь прирост на XXZ.

\subsection{Что из этого следует для выводов}

\begin{center}\small
\begin{tabular}{@{}p{9.5cm}p{4.5cm}@{}}
\toprule
Вывод & Насколько надёжен \\
\midrule
Физический prior необходим на 1 кубите & \textbf{Твёрдо} — разрыв огромный \\
Формула точна при постоянной $\gamma$ (0.9995) & \textbf{Твёрдо} — аналитика \\
Раннее предупреждение работает с 5\,\% траектории & \textbf{Твёрдо} \\
Прямое предсказание бьёт обратную задачу & \textbf{Твёрдо} — $R^2<0$ везде \\
Transformer лучше на 2 кубитах & \textbf{Твёрдо} — 0.897 против 0.636 \\
\code{physics\_lstm} хуже \code{lstm\_only} на XXZ & \textbf{Не разрешается} — $0.599\pm0.120$ против $0.633\pm0.060$ \\
Оптимальное $\tau=0.5$ для XXZ & \textbf{Отозвано} — на 3 сидах эффекта нет \\
\code{physics\_lstm} не масштабируется данными & \textbf{Опровергнуто} — знак меняется по сидам \\
E11a меняет $R^2$ на AUROC & \textbf{Опровергнуто} — окно 20 лучше по обеим \\
TFIM-потолок структурен & \textbf{Теория устояла}, но провал у перехода — нет \\
Провал предсказуемости у $J\approx h$ & \textbf{Опровергнуто} — минимум кочует по сидам \\
\bottomrule
\end{tabular}
\end{center}

\textbf{Что из этого уже сделано:} абляция E6a и $\tau$-скан E7b
реплицированы на трёх сидах при полном размере — именно это перевело две
строки выше из «хрупко» в «не разрешается» и «отозвано». Попутно нашлись и
закрыты три дефекта воспроизводимости (несеяный RNG обучения, несеяное
начальное квантовое состояние, обход засеянного пути самими скриптами), так
что повторные прогоны теперь побитово совпадают. Оставшиеся строки
«гипотеза» той же проверки ещё не проходили.

% ============================================================
\section{Что за пределами работ}

\begin{itemize}
  \item Всё на \textbf{симуляции}. Проверки на реальном железе не было: там
        $\gamma(t)$ не параметризуется полиномами Бернштейна, а ошибки
        считывания коррелированы.
  \item $N=2$ — конечноразмерные эффекты, а не настоящий фазовый переход.
  \item Обе статьи \textbf{не вычитаны человеком} перед подачей; аффилиация
        автора — плейсхолдер «Independent Researcher».
  \item Не сделано, хотя напрашивается: явные признаки запутанности
        (конкурренция, негативность); плотный скан $J$ около критической
        области; расширение на 3–4 кубита.
\end{itemize}

% ============================================================
\section{Где что лежит}

\begin{center}
\begin{tabular}{@{}p{6cm}p{8cm}@{}}
\toprule
Что & Где \\
\midrule
Статья 1 & \code{paper/paper\_1/main.tex}, \code{main.pdf} \\
Статья 2 & \code{paper/paper\_2/paper2.tex}, \code{paper2.pdf} \\
Рисунки & \code{paper/paper\_*/figures/}, код — \code{notebooks/figures.ipynb} \\
Результаты экспериментов & \code{experiments/results/*.csv} \\
Этот документ & \code{paper/papers\_explained\_ru.tex} \\
Markdown-версия & \code{docks/main/PAPERS\_EXPLAINED\_RU.md} \\
Сборка PDF & \code{make pdf} $\cdot$ \code{make pdf-1} $\cdot$ \code{make pdf-2}
  $\cdot$ \code{make pdf-ru} \\
\bottomrule
\end{tabular}
\end{center}

\end{document}
</content>
